Let $A\subseteq \mathbb{N}$ be infinite. Assume $A$ satisfies the following property:
\begin{quote}
(\*) There do not exist \emph{distinct} $a,b,c\in A$ such that $b>a$, $c>a$, and $a\mid (b+c)$.
\end{quote}
Questions:
\begin{enumerate}
\item[(Q1)] Does there exist such an $A$ with
\[\liminf_{N\to\infty} \frac{|A\cap\{1,\dots,N\}|}{N^{1/2}}>0?\]
\item[(Q2)] Does there exist an absolute constant $c>0$ such that for every such $A$ there are infinitely many $N$ with
\[|A\cap\{1,\dots,N\}|<N^{1-c}?\]
\item[(Q3)] Is it true that for every such $A$ one has
\[\sum_{n\in A}\frac{1}{n}<\infty?\]
\end{enumerate}

Stress points:
\begin{itemize}
\item The condition only forbids divisibility by an element $a$ that is \emph{smaller} than the other two elements.
\item For finite truncations, the property can hold vacuously for large elements because there may not be two larger elements to pair with a given $a$.
\end{itemize}

2)
